% Generated draft: requires mathematical and source review.
\documentclass[11pt]{article}
\usepackage[margin=25mm]{geometry}
\usepackage{fontspec}
\setmainfont{Latin Modern Roman}
\usepackage{amsmath,amssymb,mathtools,unicode-math}
\setmathfont{Latin Modern Math}
\usepackage{enumitem,needspace,xurl,hyperref}
\setlength{\parindent}{0pt}
\setlength{\parskip}{5pt}
\setlength{\emergencystretch}{3em}
\newcommand{\N}{\mathbb{N}}
\newcommand{\Z}{\mathbb{Z}}
\newcommand{\Q}{\mathbb{Q}}
\newcommand{\R}{\mathbb{R}}
\newcommand{\C}{\mathbb{C}}
\newcommand{\F}{\mathbb{F}}
\newcommand{\A}{\mathbb{A}}
\newcommand{\PP}{\mathbb{P}}
\newcommand{\E}{\mathbb{E}}
\newcommand{\eps}{\varepsilon}
\newcommand{\dd}{\,\mathrm{d}}
\begin{document}
\section*{UnsolvedMath ID: 30004217}
\textbf{Source problem code:} catalogue entry
\subsection{Definitions and mathematical statement}
The supplied record asks for a classification of triples \((n,d,k)\), but does not define the role or range of \(k\), the admissible subdivisions, or the extension relation. The following standard geometric model makes a possible interpretation explicit; its identification with the source question requires mathematical review.

Let \(d\geq 1\), \(n\geq d+1\), and \(t_1<\cdots<t_n\) be real numbers. Set
\[
 a_i=(t_i,t_i^2,\ldots,t_i^d),\qquad
 C_{n,d}=\operatorname{conv}\{a_1,\ldots,a_n\}.
\]
Write \(\widehat a_i=(a_i,1)\), and define the linear map
\[
 \pi:\R^n\longrightarrow\R^{d+1},\qquad
 \pi(x)=\sum_{i=1}^n x_i\widehat a_i.
\]
In the homogenized convention, the associated cyclic zonotope is
\[
 Z=\pi([0,1]^n)=\sum_{i=1}^n[0,\widehat a_i].
\]
For an integer \(0\leq k\leq n\), define the hypersimplex and its image by
\[
 \Delta(n,k)=\left\{x\in[0,1]^n:\sum_{i=1}^n x_i=k\right\},
 \qquad P_k=\pi(\Delta(n,k)).
\]
Thus \(P_k\) is the section of \(Z\) by the hyperplane \(H_k=\{(y,s):s=k\}\), and \(P_1=C_{n,d}\times\{1\}\). In particular, for general \(k\), the object being subdivided in this model is \(P_k\), rather than literally \(C_{n,d}\).

For this possible interpretation, call a subdivision hypersimplicial if it is a finite polytopal subdivision of \(P_k\) whose cells are images under \(\pi\) of faces of \(\Delta(n,k)\). Here a polytopal subdivision covers its underlying polytope, contains every face of each cell, and has intersections that are common faces. Model a zonotopal tiling as a polytopal subdivision of \(Z\) whose cells are images under \(\pi\) of faces of \([0,1]^n\). Its restriction to \(H_k\) consists of the nonempty intersections of its cells with \(H_k\), with repeated cells identified.

Under precisely these conventions, the classification question would be to determine
\[
 \left\{(n,d,k):
 \begin{array}{l}
 \text{every hypersimplicial subdivision }\mathcal S\text{ of }P_k\\
 \text{is the restriction to }H_k\text{ of some zonotopal tiling }\mathcal T\text{ of }Z
 \end{array}\right\}.
\]
The displayed parameter ranges and definitions specify this possible reading only. The supplied record does not establish that it intends these ranges, allows arbitrary rather than fine subdivisions or tilings, uses equality of the restricted subdivision rather than refinement, or imposes no additional conditions on projected face labels. It also does not specify whether the property is required for a fixed realization \((t_i)\) or for every realization. Consequently, this is a conditional formulation, not a fully recovered definitive statement of the source problem.
\subsection{Short English statement}
Classify the parameters for which every hypersimplicial subdivision can be obtained as a level section of a tiling of the associated cyclic zonotope. The supplied statement leaves the meaning of \(k\) and several subdivision and lifting conventions unspecified; the mathematical formulation above is explicitly conditional. The record labels the problem partially solved, but its dated literature assessment has not been independently verified.
\subsection{Sources}
\begin{itemize}
\item Local UnsolvedMath registry, permanent ID \(30004217\), source problem code \texttt{catalogue entry}, title: ``Lifting Hypersimplicial Subdivisions to Zonotopal Tilings.'' The supplied record attributes the collection to Oberwolfach Reports workshop discussions, records proposed year 2019, and supplies no named proposer.
\item Supplied source citation: \emph{Geometric, Algebraic, and Topological Combinatorics} (2020), Oberwolfach Reports, pp.\ 2395--2472, DOI \texttt{10.4171/owr/2019/39}. \url{https://doi.org/10.4171/owr/2019/39}. No problem-specific page or passage is supplied. The citation year 2020 and the proposed year 2019 are retained as recorded.
\item Supplied literature reference: Olarte and Santos, ``hypersimplicial subdivisions.'' \url{https://doi.org/10.1007/s00029-021-00722-6}.
\item Supplied literature reference: Olarte and Santos, ``primary preprint.'' \url{https://arxiv.org/abs/1906.05764}.
\item Literature assessment dated exactly \texttt{2026-08-22}: supplied dated curation, not independently verified. The record labels the problem \texttt{partially\_solved} and states that non-liftability is proved for every odd cyclic-polytope dimension, while the parameter classification and remaining even-dimensional cases are incomplete. The record itself requests independent verification of time-sensitive status and cited claims. It gives no parameter-specific theorem supporting that assessment.
\item The supplied statement status is \texttt{corrected\_verified}; its accompanying explanation says that notation, typography, and sentence structure were repaired while preserving mathematical content. This is recorded editorial provenance, not independent mathematical verification of this draft.
\end{itemize}
\end{document}
